% 固定两个输入，比较阈值假设空间与区间假设空间能够实现的标签安排。
\begin{tikzpicture}[x=1mm,y=1mm,font=\footnotesize]
  \fill[BookPaper] (0,0) rectangle (108,83);
  \node[font=\kaishu\footnotesize,text=BookInk] at (54,77)
    {固定输入 $C=\{x_1,x_2\}$，其中 $x_1<x_2$};
  \node[font=\kaishu\footnotesize,text=BookMuted] at (29,66) {目标标签};
  \node[font=\kaishu\footnotesize,text=BookMuted] at (68,66) {阈值假设};
  \node[font=\kaishu\footnotesize,text=BookMuted] at (94,66) {区间假设};
  \draw[BookRule] (4,61) -- (104,61);
  \draw[BookRule] (54,11) -- (54,72);
  \draw[BookRule] (82,11) -- (82,72);
  \foreach \yy/\a/\b in {54/0/0,43/0/1,32/1/0,21/1/1}{
    \draw[BookMuted,thin] (12,\yy) -- (36,\yy);
    \ifnum\a=1
      \fill[FigureInput] (18,\yy) circle (1.7);
    \else
      \draw[FigureLoss,thick,fill=BookPaper] (18,\yy) circle (1.7);
    \fi
    \ifnum\b=1
      \fill[FigureInput] (30,\yy) circle (1.7);
    \else
      \draw[FigureLoss,thick,fill=BookPaper] (30,\yy) circle (1.7);
    \fi
    \node[text=BookInk] at (45,\yy) {$(\a,\b)$};
    \node[text=FigureModel] at (94,\yy) {可实现};
  }
  \node[text=FigureModel] at (68,54) {可实现};
  \node[text=FigureModel] at (68,43) {可实现};
  \node[text=FigureLoss] at (68,32) {无法实现};
  \node[text=FigureModel] at (68,21) {可实现};
  \draw[BookRule] (4,14) -- (104,14);
  \node[font=\kaishu\footnotesize,text=FigureLoss] at (68,6) {未打散 $C$};
  \node[font=\kaishu\footnotesize,text=FigureModel] at (94,6) {打散 $C$};
\end{tikzpicture}
